\documentclass[10pt,a4paper]{paper}
\usepackage[utf8]{inputenc}
\usepackage[english]{babel}
\usepackage{tikz}
\usepackage{amsmath}
\usepackage{amsfonts}
\usepackage{amssymb}
\usepackage{bm}
\usepackage{geometry}
\usepackage{booktabs}
\usepackage{graphicx}
\usepackage{subcaption}

\geometry{margin=1in}

% Paragraphs
\setlength{\parindent}{4pt}
%\setlength{\parskip}{1\baselineskip}

% Notation
\newcommand{\chin}[2]{\chi_{(#1,#2)}}          % MPPI key chi_(n,m)
\newcommand{\dchi}{\delta\chi}                  % pump induced change of chi
\newcommand{\deps}{\Delta\varepsilon}           % transient Delta eps = 4 pi dchi
\newcommand{\EP}{E_{\mathrm{P}}}                % pump field
\newcommand{\Ep}{E_{\mathrm{p}}}                % probe field
\newcommand{\wP}{\omega_{\mathrm{P}}}           % pump frequency
\newcommand{\tP}{t_{\mathrm{P}}}                % pump maximum
\newcommand{\tp}{t_{\mathrm{p}}}                % probe time
\newcommand{\fenv}{f}                           % pump envelope


\title{Transient absorption of LiF from the equilibrium non-linear susceptibilities}
\author{Marco D'Alessandro, Davide Sangalli}
\date{\today}

\begin{document}
\maketitle

\section{Pump-probe response from the non-linear susceptibilities}
\label{sec:theory}

This section collects the relations between the polarization and the equilibrium non-linear susceptibilities that
we use to compute the transient response of a pump-probe experiment, and the approximations behind them. The
notation of the codes and the details of the computations are given in the next sections.

\subsection{Polarization and non-linear susceptibilities}
\label{sec:theory_chi}

Fields and polarization are written as superpositions of complex exponentials~\cite{Boyd},
\begin{equation}
\bm{E}(t) = \int d\omega\,\bm{E}(\omega)\,e^{-i\omega t} , \qquad
\bm{P}(t) = \int d\omega\,\bm{P}(\omega)\,e^{-i\omega t} ,
\label{eq:fourier}
\end{equation}
so that $\bm{E}(\omega) = (2\pi)^{-1}\int dt\,\bm{E}(t)e^{i\omega t}$ and $\bm{E}(-\omega) = \bm{E}(\omega)^*$ for a
real field. Formulas are in atomic units ($\hbar = 1$); numbers are given in eV and fs ($\hbar = 0.6582$~eV\,fs).
The polarization is expanded in powers of the field, $\bm{P} = \bm{P}^{(1)}+\bm{P}^{(2)}+\bm{P}^{(3)}+\dots$, with
\begin{align}
P^{(1)}_i(\omega) &= \sum_j\chi^{(1)}_{ij}(\omega)\,E_j(\omega) , \label{eq:P1}\\
P^{(2)}_i(\omega) &= \sum_{jk}\int d\omega_1 d\omega_2\,\delta(\omega-\omega_1-\omega_2)\,
\chi^{(2)}_{ijk}(\omega;\omega_1,\omega_2)\,E_j(\omega_1)E_k(\omega_2) , \label{eq:P2}\\
P^{(3)}_i(\omega) &= \sum_{jkl}\int d\omega_1 d\omega_2 d\omega_3\,\delta(\omega-\omega_1-\omega_2-\omega_3)
\nonumber\\
&\quad\times\chi^{(3)}_{ijkl}(\omega;\omega_1,\omega_2,\omega_3)\,E_j(\omega_1)E_k(\omega_2)E_l(\omega_3) .
\label{eq:P3}
\end{align}
The first frequency argument of $\chi^{(n)}$ is the output frequency, equal to the sum of the input ones. Since the
integration variables and the indices are dummy, $\chi^{(n)}$ can be taken symmetric under the simultaneous
permutation of the input frequencies and of the corresponding Cartesian indices (intrinsic permutation symmetry),
e.g.\ $\chi^{(2)}_{ijk}(\omega;\omega_1,\omega_2) = \chi^{(2)}_{ikj}(\omega;\omega_2,\omega_1)$.

\subsection{Pump and probe fields}
\label{sec:theory_fields}

In a pump-probe setup the field is the sum of a weak probe and of a pump,
$\bm{E}(t) = \bm{E}_{\mathrm{p}}(t)+\bm{E}_{\mathrm{P}}(t)$. We take an impulsive (delta) probe at the time $\tp$,
\begin{equation}
\bm{E}_{\mathrm{p}}(t) = E_{0\mathrm{p}}\,\bm{e}_{\mathrm{p}}\,\delta(t-\tp) , \qquad
\bm{E}_{\mathrm{p}}(\omega) = \frac{E_{0\mathrm{p}}}{2\pi}\,\bm{e}_{\mathrm{p}}\,e^{i\omega\tp} ,
\label{eq:delta_probe}
\end{equation}
so that all the probe frequencies have the same amplitude and a phase fixed by $\tp$. In our simulations probe and
pump are along the same direction, $\bm{e}_{\mathrm{p}} = \bm{e}_{\mathrm{P}} = \hat{\bm{x}}$, and we need only $P_x$.
From now on we drop the vector notation and the indices: $E_{\mathrm{p}}$, $\EP$, $P$ are the $x$ components and
$\chi^{(n)}$ the $xx\dots x$ components.

Inserting the sum of the two fields in Eqs.~\eqref{eq:P1}-\eqref{eq:P3}, the polarization is the one with the pump
only, $P_{\mathrm{P}}$ (all orders in the pump), plus the terms that contain the probe. We keep the terms linear in
the probe and up to quadratic in the pump. In $\chi^{(n)}$ the probe can take any of the $n$ input places, which give
the same contribution by permutation symmetry: this gives the factors 2 and 3 below. Calling $\omega-\Omega$ the
frequency of the probe and $\Omega$ the total frequency given by the pump,
\begin{align}
P(\omega) &= P_{\mathrm{P}}(\omega) + \chi^{(1)}(\omega)\,E_{\mathrm{p}}(\omega)
+ 2\int d\Omega\,\chi^{(2)}(\omega;\omega-\Omega,\Omega)\,E_{\mathrm{p}}(\omega-\Omega)\,\EP(\Omega) \nonumber\\
&+ 3\int d\Omega\int d\omega_1\,\chi^{(3)}(\omega;\omega-\Omega,\omega_1,\Omega-\omega_1)\,
E_{\mathrm{p}}(\omega-\Omega)\,\EP(\omega_1)\,\EP(\Omega-\omega_1) \nonumber\\
&+ O(E_{\mathrm{p}}^2,\,E_{\mathrm{p}}\EP^3) .
\label{eq:P_pp}
\end{align}
The terms of order $E_{\mathrm{p}}\EP^n$ describe the probe photon at $\omega-\Omega$ that leaves at $\omega$ after
the exchange of $n$ pump photons; for the delta probe~\eqref{eq:delta_probe}
$E_{\mathrm{p}}(\omega-\Omega) = E_{\mathrm{p}}(\omega)\,e^{-i\Omega\tp}$.

\subsection{Non-equilibrium response}
\label{sec:theory_neq}

The transient response is defined as the polarization induced by the probe in presence of the pump, divided by the
probe field,
\begin{equation}
\chineq[\EP](\omega,\tp) = \frac{\PPp(\omega)-\PPo(\omega)}{E_{\mathrm{p}}(\omega)} ,
\label{eq:chineq}
\end{equation}
where $\PPp$ is the polarization with both the fields and $\PPo$ the one with the pump only. The difference removes
$P_{\mathrm{P}}$ and, by Eq.~\eqref{eq:P_pp}, $\chineq = \chi^{(1)}+\dchi$ with the pump induced change
$\dchi = \dchi^{(2)}+\dchi^{(3)}+\dots$,
\begin{align}
\dchi^{(2)}(\omega,\tp) &= 2\int d\Omega\,e^{-i\Omega\tp}\,\chi^{(2)}(\omega;\omega-\Omega,\Omega)\,\EP(\Omega) ,
\label{eq:dchi2}\\
\dchi^{(3)}(\omega,\tp) &= 3\int d\Omega\,e^{-i\Omega\tp}\int d\omega_1\,
\chi^{(3)}(\omega;\omega-\Omega,\omega_1,\Omega-\omega_1)\,\EP(\omega_1)\EP(\Omega-\omega_1) ,
\label{eq:dchi3}
\end{align}
and in general, for the term with $n$ pump fields,
\begin{align}
\dchi^{(n+1)}(\omega,\tp) &= (n+1)\int d\Omega\,e^{-i\Omega\tp}\int d\omega_1\dots d\omega_n\,
\delta\Big(\Omega-\sum_{i=1}^n\omega_i\Big) \nonumber\\
&\quad\times\chi^{(n+1)}(\omega;\omega-\Omega,\omega_1,\dots,\omega_n)\prod_{i=1}^n\EP(\omega_i) .
\label{eq:dchin}
\end{align}
No assumption is made on the pump. The frequency $\Omega$ given by the pump appears as a phase $e^{-i\Omega\tp}$, which
is how $\dchi$ depends on the probe time. If $\chi^{(3)}$ did not depend on $\Omega$, the integrals in
Eq.~\eqref{eq:dchi3} would give, by the convolution theorem,
\begin{equation}
\int d\Omega\,e^{-i\Omega t}\int d\omega_1\,\EP(\omega_1)\EP(\Omega-\omega_1) = \EP(t)^2 ,
\label{eq:convolution}
\end{equation}
i.e.\ the response to the pump field at the probe time (and similarly $\EP(\tp)$ in $\dchi^{(2)}$).

\subsection{Monochromatic pump}
\label{sec:theory_mono}

We consider a monochromatic pump of amplitude $\EOP$, frequency $\wP$ and phase fixed by the time $\tP$,
\begin{equation}
\EP(t) = \EOP\sin(\wP(t-\tP)) = \EP(\wP)\,e^{-i\wP t}+\EP(-\wP)\,e^{i\wP t} ,
\label{eq:sine_pump}
\end{equation}
with the complex amplitudes $\EP(\pm\wP) = \pm(i\EOP/2)\,e^{\pm i\wP\tP}$. The pump frequencies in Eqs.~\eqref{eq:dchi2}-\eqref{eq:dchi3} take only the values $\pm\wP$, and the integrals become
sums over them. In the second order term the pump gives $\Omega = \pm\wP$,
\begin{equation}
\dchi^{(2)}(\omega,\tp) = \chi^{(2)}_+(\omega)\,\EP(\wP)\,e^{-i\wP\tp} + \chi^{(2)}_-(\omega)\,\EP(-\wP)\,e^{i\wP\tp} ,
\label{eq:dchi2_mono}
\end{equation}
with
\begin{equation}
\chi^{(2)}_\pm(\omega) = 2\,\chi^{(2)}(\omega;\omega\mp\wP,\pm\wP) .
\label{eq:chi2_mono}
\end{equation}
In the third order term it gives $\Omega = 0$ (in two ways, $\omega_1 = \pm\wP$) and $\Omega = \pm2\wP$,
\begin{equation}
\dchi^{(3)}(\omega,\tp) = \chiz(\omega)\,|\EP(\wP)|^2 + \chip(\omega)\,\EP(\wP)^2\,e^{-2i\wP\tp}
+ \chim(\omega)\,\EP(-\wP)^2\,e^{2i\wP\tp} ,
\label{eq:dchi3_mono}
\end{equation}
with
\begin{equation}
\chiz(\omega) = 6\,\chi^{(3)}(\omega;\omega,\wP,-\wP) , \qquad
\chi^{(3)}_{\pm}(\omega) = 3\,\chi^{(3)}(\omega;\omega\mp2\wP,\pm\wP,\pm\wP) .
\label{eq:chi3_mono}
\end{equation}
$\chiz$ is the change of the response at the probe frequency (one pump photon absorbed and one emitted: Stark shift,
bleaching, \dots). In $\chi^{(2)}_\pm$ and $\chi^{(3)}_\pm$ the probe component at $\omega\mp\wP$ or $\omega\mp2\wP$ is
converted to $\omega$ by one or two pump photons. These terms contribute to $\dchi$ at $\omega$ because the delta probe
contains all the frequencies with a fixed relative phase: they interfere with the probe at $\omega$ and oscillate with
$\tp$.

Inserting the amplitudes~\eqref{eq:sine_pump}, $|\EP(\wP)|^2 = \EOP^2/4$, and the phases combine into the delay
$\tau = \tp-\tP$:
\begin{align}
\dchi(\omega,\tau) &= \frac{i\EOP}{2}\left[\chi^{(2)}_+(\omega)\,e^{-i\wP\tau} - \chi^{(2)}_-(\omega)\,e^{i\wP\tau}\right]
\nonumber\\
&+ \frac{\EOP^2}{4}\left[\chiz(\omega) - \chip(\omega)\,e^{-2i\wP\tau} - \chim(\omega)\,e^{2i\wP\tau}\right]
+ O(\EOP^3) .
\label{eq:dchi_mono}
\end{align}
Only the delay enters, i.e.\ the phase of the pump at the probe time. The term linear in the pump oscillates with
period $2\pi/\wP$ (2.67~fs at $\wP = 1.55$~eV) and changes sign with the pump, $\EOP\to-\EOP$; it vanishes in systems
with inversion symmetry, where all the even order susceptibilities are zero. In the quadratic term the $\chiz$ part does
not depend on $\tau$ and the $\chi^{(3)}_\pm$ parts oscillate with period $\pi/\wP$ (1.33~fs); the minus signs come
from $(\pm i)^2$, i.e.\ from the sine shape of the pump (a cosine pump gives $+$).

\subsection{Pulsed pump}
\label{sec:theory_pulse}

We now consider a pump pulse with a Gaussian envelope,
\begin{equation}
\EP(t) = \EOP\,\fenv(t-\tP)\sin(\wP(t-\tP)) = a(t)\,e^{-i\wP t}+a^*(t)\,e^{i\wP t} , \qquad
\fenv(t) = e^{-t^2/2\sigma^2} ,
\label{eq:pump_pulse}
\end{equation}
with the slowly varying amplitude $a(t) = (i\EOP/2)\,\fenv(t-\tP)\,e^{i\wP\tP}$ ($\sigma = 3.2$~fs in the RT runs,
Sec.~\ref{sec:rt}). The spectrum of the pulse has two peaks, $\EP(\omega) = E^+(\omega)+E^-(\omega)$, with
\begin{equation}
E^\pm(\omega) = \pm\frac{i\EOP}{2}\,\frac{\sigma}{\sqrt{2\pi}}\,e^{i\omega\tP}\,e^{-(\omega\mp\wP)^2\sigma^2/2} ,
\label{eq:pump_spectrum}
\end{equation}
centered at $\pm\wP$ with width $1/\sigma$ ($\hbar/\sigma\simeq0.2$~eV): $E^+$ and $E^-$ are the spectra of the two
terms $a(t)e^{-i\wP t}$ and $a^*(t)e^{i\wP t}$ of Eq.~\eqref{eq:pump_pulse}.

\paragraph{Second order.} With $\EP = E^++E^-$, Eq.~\eqref{eq:dchi2} is the sum of two terms,
\begin{align}
\dchi^{(2)}(\omega,\tp) &= 2\int d\Omega\,e^{-i\Omega\tp}\,\chi^{(2)}(\omega;\omega-\Omega,\Omega)\,E^+(\Omega)
\nonumber\\
&+ 2\int d\Omega\,e^{-i\Omega\tp}\,\chi^{(2)}(\omega;\omega-\Omega,\Omega)\,E^-(\Omega) ,
\label{eq:dchi2_pulse}
\end{align}
in which the Gaussians $E^\pm$ restrict the pump frequency $\Omega$ to a range of width $1/\sigma$ around $+\wP$ or
$-\wP$. No approximation is made. For $\sigma\to\infty$ the normalized Gaussians become $\delta$-functions,
$(\sigma/\sqrt{2\pi})\,e^{-(\Omega\mp\wP)^2\sigma^2/2}\to\delta(\Omega\mp\wP)$, and
\begin{equation*}
E^\pm(\Omega)\;\xrightarrow{\;\sigma\to\infty\;}\;\pm\frac{i\EOP}{2}\,e^{\pm i\wP\tP}\,\delta(\Omega\mp\wP)
= \EP(\pm\wP)\,\delta(\Omega\mp\wP) ,
\end{equation*}
with the amplitudes $\EP(\pm\wP)$ of the monochromatic pump~\eqref{eq:sine_pump}: the integrals select
$\Omega = \pm\wP$ and Eq.~\eqref{eq:dchi2_pulse} becomes Eq.~\eqref{eq:dchi2_mono}, with
$2\chi^{(2)}(\omega;\omega\mp\wP,\pm\wP) = \chi^{(2)}_\pm(\omega)$, Eq.~\eqref{eq:chi2_mono}. For a finite $\sigma$
these are the values of the susceptibilities at the centers of the two ranges.

\paragraph{Third order.} The two pump fields in Eq.~\eqref{eq:dchi3} give four products
$E^a(\omega_1)E^b(\Omega-\omega_1)$, $a,b = \pm$, each different from zero only for $\omega_1$ close to $a\wP$ and
$\Omega$ close to $(a+b)\wP$: $E^+E^-$ and $E^-E^+$ with $\Omega\simeq0$, $E^+E^+$ with $\Omega\simeq2\wP$ and $E^-E^-$
with $\Omega\simeq-2\wP$. At fixed $\Omega$ the integral over $\omega_1$ runs over the ways $\Omega$ is shared between
the two pump photons, with $\omega_1$ within $1/\sigma$ of $\pm\wP$. In $\chi^{(3)}(\omega;\omega-\Omega,\omega_1,
\Omega-\omega_1)$ a change of $\omega_1$ at fixed $\Omega$ changes only the two pump arguments, which stay close to
$\pm\wP$, far below the gap: there $\chi^{(3)}$ changes on a scale of eV, much larger than $1/\sigma$. So $\chi^{(3)}$
can be taken out of the integral over $\omega_1$, at the center of the peak $\omega_1 = \pm\wP$ (the second pump
argument, $\Omega-\omega_1$, follows from the sum of the frequencies). A change of $\Omega$ also changes the probe
argument $\omega-\Omega$, which is resonant: this dependence is kept. The two cross terms give the same contribution
(they exchange the two pump photons) and
\begin{align}
\dchi^{(3)}(\omega,\tp) &= 6\int d\Omega\,e^{-i\Omega\tp}\,\chi^{(3)}(\omega;\omega-\Omega,\wP,\Omega-\wP)\,S_0(\Omega)
\nonumber\\
&+ 3\int d\Omega\,e^{-i\Omega\tp}\,\chi^{(3)}(\omega;\omega-\Omega,\wP,\Omega-\wP)\,S_{+2}(\Omega) \nonumber\\
&+ 3\int d\Omega\,e^{-i\Omega\tp}\,\chi^{(3)}(\omega;\omega-\Omega,-\wP,\Omega+\wP)\,S_{-2}(\Omega) ,
\label{eq:dchi3_pulse}
\end{align}
where
\begin{equation}
S_0(\Omega) = \int d\omega_1\,E^+(\omega_1)\,E^-(\Omega-\omega_1) , \qquad
S_{\pm2}(\Omega) = \int d\omega_1\,E^\pm(\omega_1)\,E^\pm(\Omega-\omega_1)
\label{eq:S3_def}
\end{equation}
are the spectra of the pairs of pump fields. With the Gaussians~\eqref{eq:pump_spectrum} the integrals over $\omega_1$ are Gaussian too,
\begin{equation}
S_0(\Omega) = \frac{\EOP^2}{4}\,\frac{\sigma}{2\sqrt{\pi}}\,e^{i\Omega\tP}\,e^{-\Omega^2\sigma^2/4} , \qquad
S_{\pm2}(\Omega) = -\frac{\EOP^2}{4}\,\frac{\sigma}{2\sqrt{\pi}}\,e^{i\Omega\tP}\,e^{-(\Omega\mp2\wP)^2\sigma^2/4} :
\label{eq:S3}
\end{equation}
they restrict $\Omega$ to a range of width $\sqrt{2}/\sigma$ around 0 and $\pm2\wP$ (the minus sign of $S_{\pm2}$ is
$(\pm i)^2$). At the centers the susceptibilities are those of the monochromatic pump,
$6\chi^{(3)}(\omega;\omega,\wP,-\wP) = \chiz(\omega)$ and
$3\chi^{(3)}(\omega;\omega\mp2\wP,\pm\wP,\pm\wP) = \chi^{(3)}_\pm(\omega)$, Eq.~\eqref{eq:chi3_mono}.

\paragraph{$\chi$ independent of $\Omega$: envelope at the probe time.} If the susceptibilities do not change over the
ranges of $\Omega$ selected by the pump (widths $1/\sigma$ and $\sqrt{2}/\sigma$), they can be taken out of the integrals at the centers of the
ranges, and the integrals left are Gaussian. With the delay $\tau = \tp-\tP$, at second order
\begin{equation}
\int d\Omega\,e^{-i\Omega\tp}\,E^\pm(\Omega) = \pm\frac{i\EOP}{2}\,\fenv(\tau)\,e^{\mp i\wP\tau} ,
\end{equation}
which are the two terms of the pump field~\eqref{eq:pump_pulse} at the probe time, and at third order
\begin{equation}
\int d\Omega\,e^{-i\Omega\tp}\,S_0(\Omega) = \frac{\EOP^2}{4}\,\fenv^2(\tau) , \qquad
\int d\Omega\,e^{-i\Omega\tp}\,S_{\pm2}(\Omega) = -\frac{\EOP^2}{4}\,\fenv^2(\tau)\,e^{\mp2i\wP\tau} ,
\end{equation}
with $\fenv^2(\tau) = e^{-\tau^2/\sigma^2}$. Eqs.~\eqref{eq:dchi2_pulse} and~\eqref{eq:dchi3_pulse} become
\begin{align}
\dchi^{(2)}(\omega,\tau) &= \frac{i\EOP}{2}\,\fenv(\tau)\left[\chi^{(2)}_+(\omega)\,e^{-i\wP\tau}
- \chi^{(2)}_-(\omega)\,e^{i\wP\tau}\right] , \nonumber\\
\dchi^{(3)}(\omega,\tau) &= \frac{\EOP^2}{4}\,\fenv^2(\tau)\left[\chiz(\omega) - \chip(\omega)\,e^{-2i\wP\tau}
- \chim(\omega)\,e^{2i\wP\tau}\right] ,
\label{eq:dchi_probetime}
\end{align}
i.e.\ the monochromatic result~\eqref{eq:dchi_mono} with the field envelope $\fenv(\tau)$ and the intensity envelope
$\fenv^2(\tau)$ at the probe time: the pump acts only through its field at the probe time.

\paragraph{Validity of the approximation.} The assumption is that the susceptibilities do not change over the range of
$\Omega$ selected by the pump, of width $\sim\hbar/\sigma$. For the pump arguments this is the condition already used
for the integral over $\omega_1$. The probe argument $\omega-\Omega$ is instead resonant: near a transition of energy
$\varepsilon$ the susceptibilities change with $\Omega$ on the scale of the width $\eta$ of the transition, as
$D(\omega-\Omega)$ with $D(\omega) = 1/(\varepsilon-\omega-i\eta)$. The condition is therefore $\eta\gg\hbar/\sigma$
(sharper structures, as an absorption edge, make it more severe).

In time this condition has a simple meaning. The delta probe starts at $\tp$ a dynamics of the system: it creates a
polarization that oscillates at the transition energies and lasts a time $\hbar/\eta$ (a line of width $\eta$
corresponds to a polarization damped as $e^{-\eta s}$, with $s = t-\tp$). The response at the frequency $\omega$ is built over all
this time, and the pump changes it as long as the polarization lives, not only at $\tp$. What enters is therefore the
pump intensity between $\tp$ and $\tp+\hbar/\eta$, and the approximation replaces it with its value at $\tp$: the pump
envelope must be constant during the dynamics started by the probe, $\hbar/\eta\ll\sigma$. This is the condition
$\eta\gg\hbar/\sigma$ seen in time.

The convolution theorem makes this statement precise. $S_0$ is the spectrum of the intensity envelope,
$S_0(\Omega) = (2\pi)^{-1}\int dt\,(\EOP^2/4)\,\fenv^2(t-\tP)\,e^{i\Omega t}$, so the term with $S_0$ in
Eq.~\eqref{eq:dchi3_pulse} is
\begin{equation}
6\int d\Omega\,e^{-i\Omega\tp}\,\chi^{(3)}(\omega;\omega-\Omega,\wP,\Omega-\wP)\,S_0(\Omega)
= \frac{\EOP^2}{4}\int ds\,K(\omega,s)\,\fenv^2(\tau+s) ,
\label{eq:kernel}
\end{equation}
with
\begin{equation*}
K(\omega,s) = \frac{6}{2\pi}\int d\Omega\,\chi^{(3)}(\omega;\omega-\Omega,\wP,\Omega-\wP)\,e^{i\Omega s} :
\end{equation*}
an average of the intensity envelope at the times $\tp+s$ after the probe, with a weight $K$ given by the dependence of
$\chi^{(3)}$ on $\Omega$ (the other terms behave in the same way). A constant $\chi^{(3)}$ gives
$K = \chiz\,\delta(s)$, i.e.\ Eq.~\eqref{eq:dchi_probetime}. The resonant dependence gives instead
\begin{equation}
\frac{1}{2\pi}\int d\Omega\,D(\omega-\Omega)\,e^{i\Omega s} = i\,e^{i(\omega-\varepsilon)s-\eta s}\,\theta(s) ,
\label{eq:g_transform}
\end{equation}
a weight that is zero before the probe (the pump cannot change the response to a probe that has not arrived yet) and
lasts $\hbar/\eta$ after it: the duration of the dynamics started by the probe.

The approximation is best at the pump maximum, where the envelope is flat, and fails first on its slopes: before the
maximum the intensity grows during the probe dynamics and Eq.~\eqref{eq:dchi_probetime} underestimates the effect of
the pump, after the maximum it overestimates it.

\paragraph{$\chi$ dependent on $\Omega$: accumulated envelope.} The susceptibilities of a monochromatic
pump give $\chi^{(3)}$ only at the centers $\Omega = 0,\pm2\wP$: the dependence on $\Omega$ requires a model. We formulate it for the third
order term. The probe is resonant: at $\tp$ it creates coherences between valence and conduction states, which
oscillate at the transition energies. The pump is far below the gap: the polarization it induces follows the field
adiabatically (virtual transitions, no memory) and its effect is to shift the transition energies while the probe
coherences live. For the delta probe the response is the transform of the response function
$R(s) = P(\tp+s)/E_{0\mathrm{p}}$,
\begin{equation}
\chi(\omega) = \int_0^\infty ds\,R(s)\,e^{i\omega s} , \qquad
R(s) = \frac{1}{2\pi}\int d\omega\,\chi(\omega)\,e^{-i\omega s} ,
\label{eq:R}
\end{equation}
and a transition of energy $\varepsilon$ and dipole $d$ contributes $R_\varepsilon(s) = i|d|^2e^{-i\varepsilon s-\eta s}$,
i.e.\ $|d|^2D(\omega)$ to $\chi^{(1)}$ (resonant term). We assume that the pump shifts the transition energy by
$\delta\varepsilon(t) = c_\varepsilon|a(t)|^2$ (dynamical Stark and Franz-Keldysh shifts). The coherence accumulates
the phase $-\int_{\tp}^{\tp+s}\delta\varepsilon\,dt'$ and, to first order in $\delta\varepsilon$,
\begin{equation}
\delta R_\varepsilon(s) = -ic_\varepsilon\,R_\varepsilon(s)\int_0^s|a(\tp+s')|^2\,ds' :
\label{eq:dR}
\end{equation}
all the times between the probe and the emission count the same. With a constant envelope, $|a|^2 = \EOP^2/4$,
$\delta R_\varepsilon(s) = (\EOP^2/4)\,R^{(3)}_\varepsilon(s)$ with $R^{(3)}_\varepsilon(s) = -ic_\varepsilon s\,R_\varepsilon(s)$,
whose transform is the contribution of the transition to $\chiz$ (a rigid shift of the line). With the pulse
\begin{align}
\delta R_\varepsilon(s) &= \frac{\EOP^2}{4}\,R^{(3)}_\varepsilon(s)\,W(\tau,s) , \nonumber\\
W(\tau,s) &= \frac{1}{s}\int_0^s\fenv^2(\tau+s')\,ds'
= \frac{\sqrt{\pi}\,\sigma}{2s}\left[\mathrm{erf}\Big(\frac{\tau+s}{\sigma}\Big)
-\mathrm{erf}\Big(\frac{\tau}{\sigma}\Big)\right] ,
\label{eq:W}
\end{align}
the mean intensity envelope between the probe and the emission. $W$ is the same for all the transitions, so the sum
gives the same relation for the total $R^{(3)}$, which is the transform~\eqref{eq:R} of $\chiz$: the single
$c_\varepsilon$, $\varepsilon$ and $d$ are not needed. The limits are the expected ones: $W\to\fenv^2(\tau)$ for
$s\ll\sigma$ (Eq.~\eqref{eq:dchi_probetime}) and $W\to1$ for $\sigma\to\infty$ (Eq.~\eqref{eq:dchi_mono}). In
frequency the model corresponds, for each transition, to
$6\chi^{(3)}(\omega;\omega-\Omega,\wP,\Omega-\wP)\propto c_\varepsilon\,D(\omega)\,D(\omega-\Omega)$: a resonant denominator at the output and one
at the input frequency of the probe, the dependence on $\Omega$ of Eq.~\eqref{eq:g_transform}.

The model holds for the term $\Omega\simeq0$, where the pump acts through $|a(t)|^2$. In the terms $\Omega\simeq\pm2\wP$
it acts through $a(t)^2e^{\mp2i\wP t}$, which oscillates with period 1.33~fs, comparable with $\hbar/\eta$, and the
Stark shift picture does not apply. We apply the same $W$ to their envelope and keep the phases $e^{\mp2i\wP\tau}$ (an
extension, not derived). The transient response is therefore
\begin{equation}
\dchi^{(3)}(\omega,\tau) = \frac{\EOP^2}{4}\int_0^\infty ds\,e^{i\omega s}\,W(\tau,s)
\left[R_0(s) - R_+(s)\,e^{-2i\wP\tau} - R_-(s)\,e^{2i\wP\tau}\right] ,
\label{eq:dchi_accumulated}
\end{equation}
where $R_0$ and $R_\pm$ are the transforms~\eqref{eq:R} of $\chiz$ and $\chi^{(3)}_\pm$. For $W = \fenv^2(\tau)$ it
reduces to the third order term of Eq.~\eqref{eq:dchi_probetime}. For comparison we also tested the local weight
$W = \fenv^2(\tau+s)$, in which the pump acts only at the emission time ($6\chi^{(3)}(\omega;\omega-\Omega,\wP,\Omega-\wP)\simeq\chiz(\omega-\Omega)$).
The accumulated weight is the one that reproduces the third order RT response at all the delays
(Sec.~\ref{sec:results}).

\section{Conventions}
\label{sec:conventions}

\section{Computation of $\chi$ with yambo\_nl}
\label{sec:yambo_nl}

\section{Numerical issues}
\label{sec:numerics}

\section{Real-time runs}
\label{sec:rt}

\section{Results}
\label{sec:results}

\section{Open points}
\label{sec:open}

\bibliographystyle{unsrt}
\bibliography{biblio}

\end{document}
